% =========================================================
% 07_determinants_inverses_cramers.tex
% Vectors and Linear Algebra
% =========================================================

% =========================================================
% SECTION DIVIDER
% =========================================================

\section{Determinants and Inverse Matrices}

% =========================================================
% DETERMINANT
% =========================================================

\begin{frame}{The Determinant}

The \textbf{determinant} is a scalar associated with a
square matrix.

For a $2\times2$ matrix

\[
A=
\begin{bmatrix}
a&b\\
c&d
\end{bmatrix},
\]

the determinant is

\[
\boxed{
\det(A)=ad-bc.
}
\]

\vspace{0.4cm}

The determinant provides information about the transformation
represented by the matrix.

\vspace{0.3cm}

It also plays a central role in determining whether a matrix
is invertible.

\end{frame}


% =========================================================
% 2x2 EXAMPLE
% =========================================================

\begin{frame}{Example: Determinant of a $2\times2$ Matrix}

Consider
\vspace{-0.1cm}
\[
A=
\begin{bmatrix}
3&2\\
1&4
\end{bmatrix}.
\]

Then

\[
\det(A)
=
(3)(4)-(2)(1).
\]

Therefore,

\[
\boxed{
\det(A)=10.
}
\]

Since

\[
\det(A)\neq0,
\]

the matrix is invertible.

\end{frame}


% =========================================================
% GEOMETRIC MEANING
% =========================================================

\begin{frame}{Geometric Meaning of the Determinant}

For a $2\times2$ matrix, the absolute value
\vspace{-0.1cm}
\[
|\det(A)|
\]

represents the factor by which the transformation changes
\textbf{area}.

If
\vspace{-0.1cm}
\[
|\det(A)|>1,
\]

areas are enlarged.

If
\vspace{-0.1cm}
\[
0<|\det(A)|<1,
\]

areas are reduced.

If
\vspace{-0.1cm}
\[
\boxed{\det(A)=0},
\]

the transformation collapses the plane into a lower-dimensional
space.

\end{frame}


% =========================================================
% 3D DETERMINANT
% =========================================================

\begin{frame}{Determinant in Three Dimensions}

For a $3\times3$ matrix, the absolute value of the determinant
represents the volume-scaling factor.

\vspace{0.4cm}

\[
\boxed{
|\det(A)|
=
\text{volume scaling factor}
}
\]

For example,

\[
\det(A)=2
\]

means that the transformation doubles volumes.

\vspace{0.4cm}

A negative determinant indicates that the transformation also
reverses orientation.

\end{frame}


% =========================================================
% PROPERTIES
% =========================================================

\begin{frame}{Properties of the Determinant}

For square matrices $A$ and $B$:

\[
\boxed{
\det(AB)=\det(A)\det(B)
}
\]

and

\[
\boxed{
\det(A^T)=\det(A).
}
\]

For the identity matrix,

\[
\boxed{
\det(I)=1.
}
\]

For a scalar $c$ and an $n\times n$ matrix,

\[
\boxed{
\det(cA)=c^n\det(A).
}
\]

\end{frame}


% =========================================================
% SINGULAR MATRIX
% =========================================================

\begin{frame}{Singular and Non-Singular Matrices}

A square matrix $A$ is called \textbf{singular} if

\[
\boxed{
\det(A)=0.
}
\]

A square matrix is \textbf{non-singular} if

\[
\boxed{
\det(A)\neq0.
}
\]

\vspace{0.4cm}

Geometrically, a singular transformation collapses
space into a lower-dimensional set.

\vspace{0.4cm}

For example, a plane may be collapsed into a line.

\end{frame}


% =========================================================
% INVERTIBILITY
% =========================================================

\begin{frame}{Determinant and Invertibility}

For a square matrix $A$,

\[
\boxed{
A^{-1}\text{ exists}
\iff
\det(A)\neq0.
}
\]

Equivalently,

\[
\boxed{
\det(A)=0
\iff
A\text{ is singular}.
}
\]

Therefore, the determinant provides a quick test for
whether a square matrix has an inverse.

\vspace{0.4cm}

This connects three ideas:

\[
\boxed{
\text{Determinant}
\quad\longleftrightarrow\quad
\text{Invertibility}
\quad\longleftrightarrow\quad
\text{Unique solution}
}
\]

\end{frame}


% =========================================================
% INVERSE MATRIX
% =========================================================

\begin{frame}{Inverse Matrix}

For a square matrix $A$, its inverse $A^{-1}$ satisfies

\[
\boxed{
AA^{-1}=A^{-1}A=I.
}
\]

Therefore,
\vspace{-0.1cm}
\[
A^{-1}
\]

undoes the transformation performed by $A$.

\vspace{0.5cm}

If
\vspace{-0.1cm}
\[
\vect{y}=A\vect{x},
\]

then

\[
\boxed{
\vect{x}=A^{-1}\vect{y}.
}
\]

Thus, the inverse transformation recovers the original vector.

\end{frame}


% =========================================================
% 2x2 INVERSE
% =========================================================

\begin{frame}{Inverse of a $2\times2$ Matrix}

For
\vspace{-0.1cm}
\[
A=
\begin{bmatrix}
a&b\\
c&d
\end{bmatrix},
\]

if
\vspace{-0.1cm}
\[
ad-bc\neq0,
\]

then
\vspace{-0.1cm}
\[
\boxed{
A^{-1}
=
\frac{1}{ad-bc}
\begin{bmatrix}
d&-b\\
-c&a
\end{bmatrix}.
}
\]

The denominator is precisely
\vspace{-0.1cm}
\[
\det(A).
\]

Therefore,
\vspace{-0.1cm}
\[
\det(A)\neq0
\]

is required for the inverse to exist.

\end{frame}


% =========================================================
% INVERSE EXAMPLE
% =========================================================

\begin{frame}{Example: Computing an Inverse}

Consider
\vspace{-0.1cm}
\[
A=
\begin{bmatrix}
2&1\\
1&1
\end{bmatrix}.
\]

First,

\[
\det(A)
=
(2)(1)-(1)(1)
=
1.
\]

Therefore,

\[
A^{-1}
=
\begin{bmatrix}
1&-1\\
-1&2
\end{bmatrix}.
\]

Indeed,

\[
AA^{-1}
=
\begin{bmatrix}
1&0\\
0&1
\end{bmatrix}
=
I.
\]

\end{frame}


% =========================================================
% INVERSE AND LINEAR SYSTEM
% =========================================================

\begin{frame}{Inverse Matrix and Linear Systems}

Consider

\[
A\vect{x}=\vect{b}.
\]

If $A$ is invertible,

\[
A^{-1}A\vect{x}
=
A^{-1}\vect{b}.
\]

Therefore,

\[
\boxed{
\vect{x}=A^{-1}\vect{b}.
}
\]

\vspace{0.4cm}

Hence, an invertible coefficient matrix guarantees a
unique solution for every right-hand side $\vect{b}$.

\end{frame}


% =========================================================
% COFACTORS
% =========================================================

\begin{frame}{Cofactors}

For larger matrices, determinants can be computed using
\textbf{cofactors}.

For an element $a_{ij}$,

\[
\boxed{
C_{ij}=(-1)^{i+j}M_{ij},
}
\]

where $M_{ij}$ is the determinant obtained after deleting
row $i$ and column $j$.

\vspace{0.4cm}

For a $3\times3$ matrix,

\[
\det(A)
=
a_{11}C_{11}
+
a_{12}C_{12}
+
a_{13}C_{13}.
\]

\vspace{0.3cm}

This is called \textbf{cofactor expansion}.

\end{frame}


% =========================================================
% ADJUGATE
% =========================================================

\begin{frame}{Adjugate and the Inverse}

The matrix of cofactors is called the \textbf{cofactor matrix}.

Its transpose is the \textbf{adjugate}:

\[
\operatorname{adj}(A)
=
C^T.
\]

For an invertible matrix,

\[
\boxed{
A^{-1}
=
\frac{1}{\det(A)}
\operatorname{adj}(A).
}
\]

\vspace{0.4cm}

This provides a general formula for the inverse of a
square matrix.

\vspace{0.3cm}

For large matrices, however, direct inverse computation is
usually not the preferred computational approach.

\end{frame}


% =========================================================
% CRAMER'S RULE
% =========================================================

\begin{frame}{Cramer's Rule}

Consider the system

\[
A\vect{x}=\vect{b},
\]

where $A$ is an invertible $n\times n$ matrix.

Cramer's rule gives each component of the solution as

\[
\boxed{
x_i
=
\frac{\det(A_i)}{\det(A)},
}
\]

where $A_i$ is obtained by replacing the $i$-th column of $A$
with $\vect{b}$.

\vspace{0.4cm}

The rule requires

\[
\boxed{
\det(A)\neq0.
}
\]

\end{frame}


% =========================================================
% CRAMER 2x2
% =========================================================

\begin{frame}{Cramer's Rule for a $2\times2$ System}

Consider
\vspace{-0.1cm}
\[
\begin{cases}
a_1x+b_1y=c_1,\\
a_2x+b_2y=c_2.
\end{cases}
\]
\vspace{-0.1cm}
The coefficient determinant is
\vspace{-0.1cm}
\[
D=
\begin{vmatrix}
a_1&b_1\\
a_2&b_2
\end{vmatrix}.
\]
\vspace{-0.1cm}
For $x$,
\vspace{-0.1cm}
\[
D_x=
\begin{vmatrix}
c_1&b_1\\
c_2&b_2
\end{vmatrix}.
\]
\vspace{-0.1cm}
For $y$,
\vspace{-0.1cm}
\[
D_y=
\begin{vmatrix}
a_1&c_1\\
a_2&c_2
\end{vmatrix}.
\]
\vspace{-0.1cm}
Therefore,
\vspace{-0.1cm}
\[
\boxed{
x=\frac{D_x}{D},
\qquad
y=\frac{D_y}{D}.
}
\]

\end{frame}


% =========================================================
% CRAMER EXAMPLE
% =========================================================

\begin{frame}{Example: Cramer's Rule}

Consider
\vspace{-0.1cm}
\[
\begin{cases}
2x+y=5,\\
x-y=1.
\end{cases}
\]

The coefficient determinant is
\vspace{-0.1cm}
\[
D=
\begin{vmatrix}
2&1\\
1&-1
\end{vmatrix}
=-3.
\]

For $x$,
\vspace{-0.1cm}
\[
D_x=
\begin{vmatrix}
5&1\\
1&-1
\end{vmatrix}
=-6.
\]

For $y$,
\vspace{-0.1cm}
\[
D_y=
\begin{vmatrix}
2&5\\
1&1
\end{vmatrix}
=-3.
\]
\vspace{-0.1cm}
Hence,
\vspace{-0.1cm}
\[
\boxed{
x=2,\qquad y=1.
}
\]

\end{frame}


% =========================================================
% CRAMER'S RULE VS GAUSSIAN ELIMINATION
% =========================================================

\begin{frame}{Cramer's Rule and Gaussian Elimination}

Both methods can solve

\[
A\vect{x}=\vect{b}.
\]

\vspace{0.4cm}

\begin{columns}[T,totalwidth=\textwidth]

\begin{column}{0.43\textwidth}

\begin{block}{Cramer's Rule}

Uses determinants:

\[
x_i=
\frac{\det(A_i)}{\det(A)}.
\]

Useful for theoretical derivations and small systems.

\end{block}

\end{column}

\begin{column}{0.43\textwidth}

\begin{block}{Gaussian Elimination}

Uses elementary row operations.

Efficient and widely used for numerical computation.

\end{block}

\end{column}

\end{columns}

\vspace{0.4cm}

For large systems, repeatedly computing determinants is
computationally expensive.

\end{frame}


% =========================================================
% KEY CONNECTION
% =========================================================

\begin{frame}{The Linear Algebra Connection}
\vspace{-0.1cm}
For a square system
\vspace{-0.2cm}
\[
A\vect{x}=\vect{b},
\]
\vspace{-0.2cm}
the following ideas are connected:
\vspace{-0.09cm}
\[
\boxed{
\begin{array}{c}
\det(A)\neq0\\
\Downarrow\\
A\text{ is invertible}\\
\Downarrow\\
A^{-1}\text{ exists}\\
\Downarrow\\
A\vect{x}=\vect{b}\text{ has a unique solution}\\
\Downarrow\\
\vect{x}=A^{-1}\vect{b}
\end{array}
}
\]
\vspace{-0.1cm}

This connection will be extended further through
\textbf{column space, null space, rank, and linear independence}.

\end{frame}